\subsection{Zweck der Discover-Phase}

Während die Identify-Phase den Problemraum aus der Vogelperspektive vermisst (\gls{deskresearch}, \gls{stakeholdermap}, \gls{designspace}), öffnet die Discover-Phase die erste Raute des \gls{doublediam} hin zur Zielgruppe selbst. Ziel ist es, die offen gebliebenen Annahmen und Wissenslücken empirisch zu prüfen und die Lebenswelt pensionierter Menschen aus erster Hand zu verstehen, also \emph{mit} der Zielgruppe zu sprechen statt \emph{über} sie.

Konkret verfolgen wir drei Erkenntnisziele: das Mobilitätsverhalten und die Entscheidungslogik zwischen \gls{miv} und \gls{oev} nachzuvollziehen, die emotionalen wie praktischen Barrieren und Motivatoren der \gls{oev}-Nutzung freizulegen sowie unsere Identify-Hypothesen zu prüfen. Die Phase ist bewusst divergent angelegt: Wir sammeln breit und unvoreingenommen, bevor in der Define-Phase verdichtet wird.

\subsection{Eingesetzte Methoden und Vorgehen}

Im Zentrum steht die qualitative Primärforschung: Wir kombinieren \gls{leitfadeninterview}s mit einem ergänzenden standardisierten Fragebogen sowie KI-generierten \emph{Synthetic Users} und werten das Material mit \gls{empathymap} und \gls{paingain} aus. Dem Profil des CAS entsprechend setzen wir gezielt KI-Werkzeuge unterstützend ein.

\subsubsection{Qualitative Leitfadeninterviews}

Als Leitmethode wählen wir das halbstandardisierte \gls{leitfadeninterview} \parencite{blume2025}, das die Innensicht der Befragten zugänglich macht. Grundregeln sind: keine Suggestiv- oder reinen Ja/Nein-Fragen, ein sparsames ``Warum?'' und stets nur eine Frage auf einmal. Die halbstandardisierte Form hält die Perspektiven vergleichbar, lässt aber Raum, unerwarteten Spuren zu folgen.

\subsubsection{Ergänzende Quellen und KI-Einsatz}

Um die qualitative Basis zu verbreitern, überführen wir den Leitfaden zusätzlich in einen standardisierten Fragebogen und verschicken ihn per E-Mail an weitere Personen der Zielgruppe. Passend zum Fokus des CAS auf \gls{designthinking} mit künstlicher Intelligenz nutzen wir generative Werkzeuge an mehreren Stellen: Claude AI als Sparringpartner beim Leitfaden, einen KI-basierten \emph{Virtual Trainer} zum Üben der Gesprächsführung sowie \emph{Synthetic Users} als KI-generierte Persona-Stellvertreter. Diese Hinweise verstehen wir bewusst als Ergänzung und nicht als Ersatz der echten Stimmen.

\subsubsection{Stichprobe, Durchführung und Synthese}

Den qualitativen Kern bilden zwei selbst geführte 30-minütige Tiefeninterviews im Rahmen der \emph{Testing Time} mit \textbf{Margrit} und \textbf{Hans Jörg}, beide pensioniert. Ergänzt um die Videointerviews der übrigen Projektteams liegen insgesamt acht Gespräche mit den beiden vor, dazu die Rückläufe aus Fragebogen und Synthetic Users. Zur Verdichtung der Rohdaten nutzen wir zwei Synthese-Werkzeuge: Die \gls{empathymap} ordnet die Aussagen in Quadranten und macht das Spannungsfeld zwischen Aussage und Verhalten sichtbar, die \gls{paingain} stellt Schmerzpunkte und Motivatoren gegenüber und bereitet die Priorisierung von Hebeln für die Develop-Phase vor.

\columnbreak
\subsection{Zentrale inhaltliche Ergebnisse}

\subsubsection{Mobilität ist Identität, nicht nur Fortbewegung}

\textbf{Das Auto steht für Freiheit und Autonomie.} Quer durch die Interviews zieht sich ein Grundmuster: ``Das Auto ist Freiheit, jederzeit einsteigen, wegfahren, ohne Fahrplan'' (Christian, Christof, Edna). Mobilität ist für die Befragten weniger eine Frage der reinen Fortbewegung als Ausdruck von Autonomie und Identität. Der \gls{oev} konkurriert damit nicht nur funktional, sondern symbolisch: Er muss aktiv ``gegen das Gewohnte'' gewählt werden, während das Auto die selbstverständliche Standardoption (``default'') bleibt. Gleichzeitig zeigt sich nach der Pensionierung eine bemerkenswerte Offenheit, ``Was kommt jetzt?'' (Peter, Christof, Hans Jörg) markiert einen biografischen Umbruch, der Raum für neue Gewohnheiten schafft.

\subsubsection{Das soziale Umfeld prägt das Verkehrsverhalten}

\textbf{Mobilität ist eine soziale Praxis.} Das unmittelbare Umfeld der Befragten ist fast ausschliesslich auto-orientiert: ``Ich kenne niemanden im Freundeskreis, der \gls{oev} fährt'' (Christian). Das erzeugt soziale Reibung und erschwert die Koordination von Ausflügen. Als Gegenkraft wirken jüngere Familienmitglieder (``Fahre doch mal mit dem Zug'') sowie Vereinskollegen, Pro Senectute und Veranstaltungskalender, die das Freizeitverhalten prägen (B.~Siegenthaler, Hans Jörg). Besonders stark ist der Wunsch nach gemeinsamem Erleben: Ausflüge mit Freundin, Gruppe oder Familie werden als sozialer Wert an sich beschrieben (Margrit, Peter).

\subsubsection{Barrieren: Wo der \gls{oev} im Alltag scheitert}

\textbf{Ticketkauf-Unsicherheit} ist die meistgenannte Hürde: ``Was, wenn ich das falsche Ticket löse?'' (P.~Büki, Margrit). Zonenfehler, eine als zu komplex empfundene App und die undurchsichtige Anschlussbillett-Logik führen bis zur Angst, ``fast als Schwarzfahrer'' dazustehen. \textbf{Das Sicherheitsgefühl} sinkt nach Einbruch der Dunkelheit, schlecht beleuchtete Bahnhöfe und ``komische Gestalten'' führen dazu, dass der \gls{oev} abends gemieden wird. \textbf{Gepäck und Transport} machen das Auto bei schweren Einkäufen, Garten- oder Werkzeugtransporten praktisch unverzichtbar. Hinzu kommen \textbf{Mitmenschen im \gls{oev}} (Lärm, Gerüche, Gedränge, fehlende Etikette) und \textbf{App-Updates}, die gelernte Bedienwege zerstören und das Vertrauen mit jeder Änderung schwächen. Die mit Abstand stärkste Barriere bleibt jedoch die \textbf{Gewohnheit}: Im Entscheidungsmoment taucht der \gls{oev} schlicht nicht als Option auf.

\subsubsection{Motivatoren: Wann der \gls{oev} überzeugt}

\textbf{Kein Parkplatzstress} ist der rationalste Treiber, ``Wenn ich nach Zürich gehe, nehme ich die S-Bahn'' (Christian, Christof), sobald das Auto zum Problem wird. Noch stärker wirken \textbf{Erlebnisfahrten}: GoldenPass, Nostalgiezug oder Dampfbahn Furka machen die Fahrt selbst zum Highlight, ``Die Fahrt ist das Ziel'' (Margrit, Hans Jörg). Für solche Erlebnisse sind die Befragten bereit zu zahlen und sorgfältig zu planen. Dazu kommen die \textbf{soziale Dimension} des gemeinsamen Unterwegsseins, die \textbf{gewonnene Zeit im Zug} (lesen, Landschaft geniessen, ohne Zeitdruck nach der Pensionierung) sowie \textbf{Natur, Wandern und Kultur} als starke Motivatoren. Bemerkenswert ist schliesslich die vorhandene \textbf{digitale Kompetenz}: Die \gls{zvv}-App wird als ``sensationell'' beschrieben (Hans Jörg, Peter), und Ausflüge werden mit Wetter-, Wanderweg- und Webcam-Apps akribisch geplant.

\columnbreak
\subsubsection{Zwei Profile im Kontrast: Hans Jörg und Margrit}

Die beiden von uns selbst geführten Interviews verdichten das Spektrum exemplarisch. \textbf{Hans Jörg fährt \gls{oev}, weil für ihn die Gewinne überwiegen}: Er erlebt die Fahrt als Teil des Ausflugs (``Bei einer schönen Strecke ist die Fahrt das Ziel''), ist digital souverän (``Die \gls{zvv}-App ist sensationell''), hat nach der Pensionierung Zeit und Musse und ist über Verein, Natur und Kultur stark an Freizeitziele gebunden, die der \gls{oev} bequem erschliesst. Die App nimmt ihm genau jene Unsicherheit, die andere abschreckt.

\textbf{Margrit fährt eher \emph{nicht} \gls{oev}, weil bei ihr die Schmerzpunkte überwiegen}: Die Ticketkauf-Unsicherheit wirkt als emotionale Hauptbarriere (``einmal stand ich fast als Schwarzfahrerin da''), und der \gls{oev} wird für sie erst attraktiv, wenn sie nicht allein unterwegs ist (``Am schönsten ist es, gemeinsam unterwegs zu sein''). Fehlt die Begleitung, überwiegen Hemmschwelle und Sicherheitsbedenken. Der direkte Vergleich zeigt: Nicht das Angebot selbst, sondern das Verhältnis von wahrgenommenen Gains und Pains entscheidet über die Nutzung.

\subsubsection{Schlüsselerkenntnis: Erlebnis schlägt Pflicht}

\textbf{Der \gls{oev} gewinnt über das Erlebnis, nicht über die Pflicht.} Die Befragten lassen sich nicht über Argumente der Vernunft oder Nachhaltigkeit aktivieren, sondern über das Versprechen eines attraktiven Freizeiterlebnisses. Wo das Auto an seine Grenzen stösst (Parkplatz, Stau) oder wo die Fahrt selbst zum Ziel wird, kippt die Entscheidung zugunsten des \gls{oev}. Diese Erkenntnis verschiebt den Fokus weg vom reinen Transportmittel hin zum Erlebnis-Zubringer.

Anhand der verdichteten Ergebnisse aus den Transkripten der Interviews wird mit Claude AI die Empathy Map und die Pain-/Gain-Map erstellt, die die Grundlage für die Define-Phase bilden.

\columnbreak
\subsection{Reflexion der Methode}

Die qualitativen \gls{leitfadeninterview}s erweisen sich als der richtige Hebel für diese Phase. Sie liefern reichhaltige, emotional aufgeladene Einblicke, die eine reine \gls{deskresearch} nicht erbringen könnte, und machen den Widerspruch zwischen Einstellung (``\gls{oev} ist mühsam'') und konkretem Verhalten (begeisterte Erlebnisfahrten) überhaupt erst sichtbar.

Gleichzeitig sind wir uns der typischen Verzerrungen bewusst, vor denen \textcite{blume2025} warnt. Dem \gls{confirmationbias}, also der Versuchung, vor allem Bestätigung für unsere Identify-Hypothesen zu suchen, begegnen wir mit bewusst offen und neutral formulierten Fragen. Dem \gls{haloeffekt} und einem \gls{ingroupbias} wirken wir entgegen, indem wir Protokolle anderer Teams einbeziehen und so unsere eigene Sichtweise relativieren.

Die Grenzen der Methode liegen in der Stichprobe: Sie ist klein, nicht repräsentativ und durch den E-Mail-Fragebogen tendenziell zu digital-affinen, aktiven Senioren verzerrt, betagte und offline lebende Personen sind unterrepräsentiert. Zudem zwingen uns technische Schwierigkeiten dazu, die Gespräche manuell statt systematisch transkribiert auszuwerten, was die Auswertungstiefe begrenzt. Auch die ergänzend befragten Synthetic Users ersetzen keine echten Menschen, da sie bestehende Annahmen reproduzieren und Verzerrungen verstärken können.

Im Rückblick würden wir die Interviews stärker durch teilnehmende \emph{Beobachtungen} ergänzen, etwa an Bahnhöfen und Ticketautomaten, um die berichtete Ticketkauf-Unsicherheit direkt im Kontext zu erfassen, und die Stichprobe gezielt um schwerer erreichbare Segmente erweitern. Für die nächste Phase liefern die verdichteten Ergebnisse aus Empathy Map und Pain-/Gain-Map die Grundlage, um Persona, \gls{pov} und How-Might-We-Fragen zu schärfen.

\FullWidthFigure{sections/figures/empathymap}{Empathy Map der frisch pensionierten Zielgruppe (65--70 Jahre) auf Basis der Leitfadeninterviews}

\FullWidthFigure{sections/figures/paingain}{Pain- und Gain-Map: Schmerzpunkte und Motivatoren der \gls{oev}-Nutzung}